\subsection{赋范空间的商空间与积空间}

命题 9. 设 E 是非离散赋值除环 K 上的赋范空间，H 是 E 的闭向量子空间。若对每个类 \(\dot{\mathbf{x}}\in\mathbf{E}/\mathbf{H}\) 令 \(||\dot{\mathbf{x}}||=\inf_{\mathbf{x}\in\dot{\mathbf{x}}}\|\mathbf{x}||\) ，则函数 \(||\dot{\mathbf{x}}||\) 是向量空间 E/H 上的范数，

且由这个范数定义的拓扑是 E 的拓扑被 H 所作的商。由第 1 号的注 2，\(d(\dot{\boldsymbol{x}}, \dot{\boldsymbol{y}}) = ||\dot{\boldsymbol{x}} - \dot{\boldsymbol{y}}||\) 是 E/H 上的不变度量，它定义 E 的拓扑被 H 所作的商。只需再证明 \(||t\dot{\boldsymbol{x}}|| = |t|.||\dot{\boldsymbol{x}}||\) ，而这由 \(||\dot{\boldsymbol{x}}||\) 的定义立即可得{[}第 IV 章，§ 5，第 7 号，公式 (23){]}。

范数 \(||\dot{\pmb{x}}||\) 也可解释如下：它是任意点 \(\pmb{x} \in \dot{\pmb{x}}\) 到子空间 H 的距离（在 E 中），因为 \(\dot{\pmb{x}}\) 的点正是诸点 \(\pmb{x} - \pmb{z}\) ，其中 \(\pmb{z}\) 取遍 H。

命题 10. 设 \((\mathbf{E}_i)_{1 \leqslant i \leqslant n}\) 是非离散赋值除环 K 上的赋范空间的有限族，\(\mathbf{E} = \prod_{i=1}^{n} \mathbf{E}_i\) 是积向量空间。若对每个 \(x = (x_i) \in \mathbf{E}\) 令 \(||x|| = \sup_{1 \leqslant i \leqslant n} ||x_i||\) ，则函数 \(||x||\) 是 E 上的范数，且它在 E 上定义的拓扑是诸 \(\mathbf{E}_i\) 的拓扑之积。

因为若 \(\pmb{x} = (\pmb{x}_i)\) 与 \(\pmb{y} = (\pmb{y}_i)\) ，则 \(\pmb{x} + \pmb{y} = (\pmb{x}_i + \pmb{y}_i)\) ，从而

\[
\begin{array}{r l} | | \boldsymbol {x} + \boldsymbol {y} | | = & \sup _ {i} | | \boldsymbol {x} _ {i} + \boldsymbol {y} _ {i} | | \leqslant \sup _ {i} (| | \boldsymbol {x} _ {i} | | + | | \boldsymbol {y} _ {i} | |) \\ & \leqslant \sup _ {i} | | \boldsymbol {x} _ {i} | | + \sup _ {i} | | \boldsymbol {y} _ {i} | | = | | \boldsymbol {x} | | + | | \boldsymbol {y} | |. \end{array}
\]

另一方面，显然 \(||tx|| = |t|.||x||\) ，且若 \(||x|| = 0\) ，则 \(||x_i|| = 0\) ，从而 \(x_i = 0\) （对 \(1 \leqslant i \leqslant n\) ），于是 \(x = 0\) ；故 \(||x||\) 是 E 上的范数。又关系 \(||x|| < a\) 等价于 \(n\) 个关系 \(||x_i|| < a\) ，因而范数 \(||x||\) 在 E 上定义积拓扑。

类似地可证函数 \(\sum_{i=1}^{n} \|x_i\|\) 与 \(\sqrt{\sum_{i=1}^{n} \|x_i\|^2}\) 是 E 上的范数；不等式 (2) 表明这三个范数都等价。

特别地，在（左或右）向量空间 \(K^{n}\) 中，若令

\[
p _ {1} (\boldsymbol {x}) = \sup _ {i} | x _ {i} |, \quad p _ {2} (\boldsymbol {x}) = \sum_ {i = 1} ^ {i = 1} | x _ {i} |, \quad p _ {3} (\boldsymbol {x}) = \sqrt {\sum_ {i = 1} ^ {n} | x _ {i} | ^ {2}}
\]

（对 \(\pmb{x} = (x_{i})_{1\leqslant i\leqslant n}\) ），则三个函数 \(p_1, p_2, p_3\) 是等价的范数，它们在 \(\mathbf{K}^n\) 上定义的拓扑即诸因子 \(\mathbf{K}\) 的拓扑之积。

